\subsection{LIE GROUP CASE}

PROPOSITION 11. Assume that k is R, C or a non-discrete complete ultrametric field. Let h be a splitting Cartan subalgebra of g.

\begin{enumerate}
\def\labelenumi{(\roman{enumi})}
\item
  \(\operatorname{Aut}(\mathfrak{g},\mathfrak{h})\) is a Lie subgroup of \(\operatorname{Aut}(\mathfrak{g})\) with Lie algebra ad h.
\item
  \(f(\mathrm{T}_{\mathrm{Q}})\) and \((q\circ f)(\mathrm{T}_{\mathrm{P}})\) are open subgroups of \(\operatorname {Aut}(\mathfrak{g},\mathfrak{h})\)
\item
  \(\mathrm{Aut}_e(\mathfrak{g})\) is an open subgroup of \(\mathrm{Aut}(\mathfrak{g})\) .
\item
  If \(k = \mathbf{R}\) or \(\mathbf{C}\) , \(\operatorname{Aut}_e(\mathfrak{g})\) is the identity component of \(\operatorname{Aut}(\mathfrak{g})\) , in other words \(\operatorname{Int}(\mathfrak{g})\) .
\end{enumerate}

By Chap. III, §3, no. 8, Cor. 2 of Prop. 29, and no. 10, Prop. 36, \(\operatorname{Aut}(\mathfrak{g}, \mathfrak{h})\) is a Lie subgroup of \(\operatorname{Aut}(\mathfrak{g})\) whose Lie algebra is the set of \(\operatorname{ad} x (x \in \mathfrak{g})\) such that \((\operatorname{ad} x) \mathfrak{h} \subset \mathfrak{h}\) , in other words \(\operatorname{ad} \mathfrak{h}\) .

Let \(\mathrm{H} \in \mathfrak{h}\) . There exists \(\varepsilon > 0\) with the following properties: for \(t \in k\) and \(|t| < \varepsilon\) , \(\exp(t\gamma(\mathrm{H}))\) is defined for all \(\gamma \in \mathrm{P}(\mathrm{R})\) , and the map \(\gamma \mapsto \exp(t\gamma(\mathrm{H}))\) is a homomorphism \(\sigma_t\) from \(\mathrm{P}(\mathrm{R})\) to \(k^*\) . For \(|t| < \varepsilon\) , \(\exp(t\operatorname{ad}\mathrm{H})\) is defined, induces the identity on \(\mathfrak{h}\) and induces on \(\mathfrak{g}^\alpha\) the homothety with ratio \(\sigma_t(\alpha)\) ; hence \(\exp t\operatorname{ad}\mathrm{H} \in (q \circ f)(\mathrm{T}_{\mathrm{P}})\) . This proves, in view of (i), that \((q \circ f)(\mathrm{T}_{\mathrm{P}})\) contains a neighbourhood of 1 in \(\operatorname{Aut}(\mathfrak{g},\mathfrak{h})\) , and consequently is an open subgroup of \(\operatorname{Aut}(\mathfrak{g},\mathfrak{h})\) . A fortiori, \(f(\mathrm{T}_{\mathrm{Q}})\) is an open subgroup of \(\operatorname{Aut}(\mathfrak{g},\mathfrak{h})\) .

For all \(\alpha \in \mathrm{R}\) , \(\exp \operatorname{ad} \mathfrak{g}^{\alpha} \subset \operatorname{Aut}_{e}(\mathfrak{g})\) . In view of (ii), \(\operatorname{Aut}_{e}(\mathfrak{g})\) contains a neighbourhood of 1 in \(\operatorname{Aut}(\mathfrak{g})\) , which proves (iii).

Assume that \(k = \mathbf{R}\) or \(\mathbf{C}\) . Then \(\mathrm{Aut}_e(\mathfrak{g})\) is contained in the identity component \(\mathrm{C}\) of \(\mathrm{Aut}(\mathfrak{g})\) (Chap. VII, §3, no. 1), and is open in \(\mathrm{Aut}(\mathfrak{g})\) by (iii). Thus \(\mathrm{Aut}_e(\mathfrak{g}) = \mathrm{C}\) . Finally, \(\mathrm{C} = \mathrm{Int}(\mathfrak{g})\) by Chap. III, §9, no. 8, Prop. 30 (i).

